\documentclass[11pt]{article}
\usepackage[margin=1in]{geometry}
\usepackage{booktabs}
\usepackage{amsmath}
\usepackage{amssymb}
\usepackage{multirow}
\usepackage{threeparttable}
\usepackage{graphicx}
\usepackage{microtype}
\usepackage{enumitem}

\title{\textbf{Statistical Significance and Runtime Analysis of Hybrid Quantum--Classical Emotion Classification}\\
\large Pairwise Comparisons of Exact Match Ratio (EMR) and Variational Circuit Complexity}
\author{}
\date{}

\begin{document}
\maketitle

\section*{1. Statistical Significance Analysis}

\subsection*{1.1 Data and Test Used}

All significance tests are computed on the \textbf{Exact Match Ratio (EMR)} reported in the main results table
of the paper for the EmoNoBa test set of
\[
N = 284 \text{ batches} \times 8 \text{ samples/batch} = 2272 \text{ samples}.
\]
Each configuration produces a single EMR $= \frac{\text{number of exactly matched samples}}{N}$.

Because the available data is the aggregate EMR per configuration (a single test-set run) rather than
per-sample prediction vectors or multiple independent training runs, we compare EMRs as proportions of the
shared test set using the \textbf{two-proportion $z$-test} (pooled variance):
\[
z = \frac{\hat{p}_1 - \hat{p}_2}{\sqrt{\hat{p}_{\text{pool}}(1 - \hat{p}_{\text{pool}})\left(\frac{1}{N} + \frac{1}{N}\right)}},
\qquad
\hat{p}_{\text{pool}} = \frac{x_1 + x_2}{2N},
\]
and report the two-sided $p$-value $p = 2\Phi(-|z|)$. A comparison is deemed \emph{statistically significant}
at the $\alpha = 0.05$ level when $p < 0.05$, and the effect is reported with its sign (which configuration
is better).

\paragraph{Important caveat.} The two-proportion $z$-test assumes the two predictions are \emph{independent}.
Here all configurations share the same test set and are strongly correlated (they often agree on the easy
samples), so this test \emph{underestimates} the correlation and may \emph{overstate} significance. For a
fully rigorous test one should instead use \emph{McNemar's test} on per-sample prediction vectors (paired
data), or a $t$-test / Mann--Whitney $U$ test across multiple independent training runs (seeds). The
aggregate single-run EMRs available here do not support those tests directly; the $z$-test below is the
most appropriate test possible with the current data and should be interpreted with this limitation in
mind.

\subsection*{1.2 IndicBERTv2 Encoder}

The classical baseline has $\text{EMR} = 0.4252$.

\subsubsection*{Classical vs.\ each quantum configuration}
\begin{table}[h!]
\centering
\small
\begin{threeparttable}
\begin{tabular}{llccccc}
\toprule
\textbf{Comparison} & \textbf{Qubits} & \textbf{VQC} & \textbf{EMR} & $z$ & $p$-value & \textbf{Sig.} \\
\midrule
Classical & 2 & SEL     & 0.4252 vs.\ 0.0238 & 32.42 & $<0.0001$ & YES \\
Classical & 4 & SEL     & 0.4252 vs.\ 0.2148 & 15.20 & $<0.0001$ & YES \\
Classical & 8 & SEL     & 0.4252 vs.\ 0.0497 & 29.74 & $<0.0001$ & YES \\
\addlinespace
Classical & 2 & Custom 1 & 0.4252 vs.\ 0.4230 & 0.15 & 0.8807 & no \\
Classical & 4 & Custom 1 & 0.4252 vs.\ 0.4155 & 0.66 & 0.5078 & no \\
Classical & 8 & Custom 1 & 0.4252 vs.\ 0.4265 & $-0.09$ & 0.9294 & no \\
\addlinespace
Classical & 2 & Custom 2 & 0.4252 vs.\ 0.1527 & 20.26 & $<0.0001$ & YES \\
Classical & 4 & Custom 2 & 0.4252 vs.\ 0.0519 & 29.52 & $<0.0001$ & YES \\
Classical & 8 & Custom 2 & 0.4252 vs.\ 0.4577 & $-2.21$ & 0.0274 & YES \\
\bottomrule
\end{tabular}
\end{threeparttable}
\caption{IndicBERTv2: classical baseline vs.\ each VQC configuration.}
\end{table}

\noindent\textbf{Reading:}
\begin{itemize}
    \item \textbf{SEL} is significantly \emph{worse} than the classical baseline at every register size ($p<0.0001$).
    \item \textbf{Custom 1} is statistically indistinguishable from the classical baseline at 2, 4, and 8 qubits
    (all $p>0.05$); the VQC neither helps nor hurts.
    \item \textbf{Custom 2} is significantly \emph{worse} than classical at 2 and 4 qubits ($p<0.0001$), but
    significantly \emph{better} at 8 qubits ($p=0.0274$) --- the only quantum configuration under
    IndicBERTv2 that beats the classical baseline.
\end{itemize}

\subsubsection*{Adjacent register-size comparisons within each VQC}
\begin{table}[h!]
\centering
\small
\begin{threeparttable}
\begin{tabular}{llcccc}
\toprule
\textbf{VQC} & \textbf{Comparison} & \textbf{EMR} & $z$ & $p$-value & \textbf{Sig.} \\
\midrule
SEL     & 2 qubits vs.\ 4 qubits  & 0.0238 vs.\ 0.2148 & $-19.86$ & $<0.0001$ & YES \\
SEL     & 4 qubits vs.\ 8 qubits  & 0.2148 vs.\ 0.0497 & 16.43 & $<0.0001$ & YES \\
\addlinespace
Custom 1 & 2 qubits vs.\ 4 qubits & 0.4230 vs.\ 0.4155 & 0.51 & 0.6084 & no \\
Custom 1 & 4 qubits vs.\ 8 qubits & 0.4155 vs.\ 0.4265 & $-0.75$ & 0.4527 & no \\
\addlinespace
Custom 2 & 2 qubits vs.\ 4 qubits & 0.1527 vs.\ 0.0519 & 11.21 & $<0.0001$ & YES \\
Custom 2 & 4 qubits vs.\ 8 qubits & 0.0519 vs.\ 0.4577 & $-31.39$ & $<0.0001$ & YES \\
\bottomrule
\end{tabular}
\end{threeparttable}
\caption{IndicBERTv2: adjacent register-size comparisons within each VQC.}
\end{table}

\noindent\textbf{Reading:}
\begin{itemize}
    \item \textbf{SEL} changes significantly between every pair of registers (non-monotonic, peaking at 4 qubits).
    \item \textbf{Custom 1} is statistically flat across registers (no significant change with qubit count).
    \item \textbf{Custom 2} changes significantly between every pair, with a dramatic improvement from 4 to 8 qubits
    (0.0519 $\to$ 0.4577).
\end{itemize}

\subsection*{1.3 BanglaBERT Encoder}

The classical baseline has $\text{EMR} = 0.0000$ (the classical model predicts no positive labels).

\subsubsection*{Classical vs.\ each quantum configuration}
\begin{table}[h!]
\centering
\small
\begin{threeparttable}
\begin{tabular}{llccccc}
\toprule
\textbf{Comparison} & \textbf{Qubits} & \textbf{VQC} & \textbf{EMR} & $z$ & $p$-value & \textbf{Sig.} \\
\midrule
Classical & 2 & SEL     & 0.0000 vs.\ 0.0264 & $-7.85$ & $<0.0001$ & YES \\
Classical & 4 & SEL     & 0.0000 vs.\ 0.3178 & $-32.53$ & $<0.0001$ & YES \\
Classical & 8 & SEL     & 0.0000 vs.\ 0.0387 & $-9.56$ & $<0.0001$ & YES \\
\addlinespace
Classical & 2 & Custom 1 & 0.0000 vs.\ 0.1505 & $-20.06$ & $<0.0001$ & YES \\
Classical & 4 & Custom 1 & 0.0000 vs.\ 0.3891 & $-38.04$ & $<0.0001$ & YES \\
Classical & 8 & Custom 1 & 0.0000 vs.\ 0.3895 & $-38.07$ & $<0.0001$ & YES \\
\addlinespace
Classical & 2 & Custom 2 & 0.0000 vs.\ 0.0000 & n/a & n/a\tnote{$\ddagger$} & n/a \\
Classical & 4 & Custom 2 & 0.0000 vs.\ 0.4071 & $-39.50$ & $<0.0001$ & YES \\
Classical & 8 & Custom 2 & 0.0000 vs.\ 0.4630 & $-44.26$ & $<0.0001$ & YES \\
\bottomrule
\end{tabular}
\begin{tablenotes}
\small
\item[$\ddagger$] Both configurations achieve $\text{EMR}=0.0000$; there is no difference to test.
\end{tablenotes}
\end{threeparttable}
\caption{BanglaBERT: classical baseline vs.\ each VQC configuration.}
\end{table}

\noindent\textbf{Reading:}
\begin{itemize}
    \item Because the classical BanglaBERT baseline collapses to zero EMR, \textbf{every} quantum configuration that
    achieves a non-zero EMR is automatically and significantly better ($p<0.0001$).
    \item This reflects the \emph{compensatory} role of the quantum branch discussed in the paper: when the classical
    embedding is weak, the quantum circuit provides a large and significant gain.
    \item The 2-qubit Custom 2 also produces zero EMR, so it cannot be distinguished from the classical baseline.
\end{itemize}

\subsubsection*{Adjacent register-size comparisons within each VQC}
\begin{table}[h!]
\centering
\small
\begin{threeparttable}
\begin{tabular}{llcccc}
\toprule
\textbf{VQC} & \textbf{Comparison} & \textbf{EMR} & $z$ & $p$-value & \textbf{Sig.} \\
\midrule
SEL     & 2 qubits vs.\ 4 qubits  & 0.0264 vs.\ 0.3178 & $-26.02$ & $<0.0001$ & YES \\
SEL     & 4 qubits vs.\ 8 qubits  & 0.3178 vs.\ 0.0387 & 24.58 & $<0.0001$ & YES \\
\addlinespace
Custom 1 & 2 qubits vs.\ 4 qubits & 0.1505 vs.\ 0.3891 & $-18.12$ & $<0.0001$ & YES \\
Custom 1 & 4 qubits vs.\ 8 qubits & 0.3891 vs.\ 0.3895 & $-0.03$ & 0.9779 & no \\
\addlinespace
Custom 2 & 2 qubits vs.\ 4 qubits & 0.0000 vs.\ 0.4071 & $-39.50$ & $<0.0001$ & YES \\
Custom 2 & 4 qubits vs.\ 8 qubits & 0.4071 vs.\ 0.4630 & $-3.80$ & 0.0001 & YES \\
\bottomrule
\end{tabular}
\end{threeparttable}
\caption{BanglaBERT: adjacent register-size comparisons within each VQC.}
\end{table}

\noindent\textbf{Reading:}
\begin{itemize}
    \item \textbf{SEL} peaks significantly at 4 qubits and degrades significantly at 8 qubits.
    \item \textbf{Custom 1} improves significantly from 2 to 4 qubits, then plateaus (4 vs.\ 8 not significant).
    \item \textbf{Custom 2} improves significantly from 4 to 8 qubits ($p=0.0001$).
\end{itemize}

\section*{2. Time Complexity and Runtime Analysis}

\subsection*{2.1 Asymptotic Complexity of the Quantum Branch}

Let $n_q$ denote the number of qubits, $L$ the number of variational blocks (with the associated
data re-upload), and $B$ the mini-batch size. The variational quantum circuit is simulated
exactly on the PennyLane \texttt{default.qubit} statevector simulator, which represents the
state as a complex vector of dimension $2^{n_q}$. Applying any single gate to the full
statevector costs $\mathcal{O}(2^{n_q})$ time.

For the Custom~2 circuit (one ansatz layer, data re-uploading performed twice), the gate count is
\begin{align*}
\text{Encoding block} &: \text{2 rotations per qubit} = 2n_q \text{ gates}, \\
\text{Ansatz block}   &: n_q \text{ single-qubit Rot} + n_q \text{ CZ} = 2n_q \text{ gates}, \\
\text{Total (2 repetitions)} &: 2\,(2n_q + 2n_q) = 8 n_q \text{ gates}.
\end{align*}
Hence a single circuit evaluation has cost
\[
T_{\text{circuit}} = \mathcal{O}\!\left(n_q \, 2^{n_q}\right),
\]
and, because \texttt{QuantumBottleneck} executes the circuit once per sample in a batch,
\[
T_{\text{batch}} = \mathcal{O}\!\left(B \, n_q \, 2^{n_q}\right).
\]

The dominant cost of a full forward pass is nevertheless the transformer encoder. For a
sequence of length $S$, hidden dimension $d$, and $L_t$ layers, the self-attention cost is
\[
T_{\text{BERT}} = \mathcal{O}\!\left(B\, L_t\, S^2\, d\right),
\]
which, for the large pretrained encoders used here (IndicBERTv2, $\approx$278M parameters),
dwarfs the quantum contribution in wall-clock time. The quantum statevector itself is
negligible in memory: $16 \cdot 2^{n_q}$ bytes, e.g.\ 4\,KB at $n_q = 8$.

\subsection*{2.2 Measured Runtime}

The variational circuit timing was instrumented live during training (``VQC TIMING''),
recording the total simulator time, the number of forward circuit calls, the per-batch time
($B = 8$), and the per-sample time. Two eight-qubit \textbf{Custom 2} configurations
(IndicBERTv2 encoder), two four-qubit \textbf{Custom 1} configurations (IndicBERTv2 and
BanglaBERT encoders), and two four-qubit \textbf{SEL} configurations (IndicBERTv2 and
BanglaBERT encoders) are reported.

\begin{table}[h!]
\centering
\small
\begin{threeparttable}
\begin{tabular}{lcc}
\toprule
\textbf{Quantity} & \textbf{\shortstack{Custom~2, IndicBERTv2\\(8q, run A)}} & \textbf{\shortstack{Custom~2, IndicBERTv2\\(8q, run B)}} \\
\midrule
Total VQC time (s)        & 13639.8 & 11729.2 \\
Total VQC time (min)      & 227.33  & 195.49  \\
Forward calls             & 48905   & 41228   \\
Per batch, $B=8$ (ms)     & 278.90  & 284.50  \\
Per sample (ms)           & 34.86   & 35.56   \\
\bottomrule
\end{tabular}
\end{threeparttable}
\caption{Measured variational-circuit runtime for the Custom~2 design (PennyLane statevector simulator, CPU).}
\end{table}

\begin{table}[h!]
\centering
\small
\begin{threeparttable}
\begin{tabular}{lcc}
\toprule
\textbf{Quantity} & \textbf{Custom~1, IndicBERTv2 (4q)} & \textbf{Custom~1, BanglaBERT (4q)} \\
\midrule
Total VQC time (s)        & 5374.5  & 5238.4  \\
Total VQC time (min)      & 89.58   & 87.31   \\
Forward calls             & 64259   & 59141   \\
Per batch, $B=8$ (ms)     & 83.64   & 88.58   \\
Per sample (ms)           & 10.45   & 11.07   \\
\bottomrule
\end{tabular}
\end{threeparttable}
\caption{Measured variational-circuit runtime for the Custom~1 design at 4 qubits (PennyLane statevector simulator, CPU).}
\end{table}

\begin{table}[h!]
\centering
\small
\begin{threeparttable}
\begin{tabular}{lcc}
\toprule
\textbf{Quantity} & \textbf{SEL, IndicBERTv2 (4q)} & \textbf{SEL, BanglaBERT (4q)} \\
\midrule
Total VQC time (s)        & 4029.1  & 4950.2  \\
Total VQC time (min)      & 67.15   & 82.50   \\
Forward calls             & 59141   & 59141   \\
Per batch, $B=8$ (ms)     & 68.13   & 83.70   \\
Per sample (ms)           & 8.52    & 10.46   \\
\bottomrule
\end{tabular}
\end{threeparttable}
\caption{Measured variational-circuit runtime for the SEL design at 4 qubits (PennyLane statevector simulator, CPU).}
\end{table}

\noindent\textbf{Reading:}
\begin{itemize}
    \item The per-sample VQC time for \textbf{Custom~2} is $\approx$35\,ms and is essentially
    identical across both eight-qubit runs of the same encoder, confirming that the quantum cost
    is dominated by the $2^{n_q}$ statevector dimension and is reproducible run-to-run.
    \item The ratio of total time to forward calls matches the per-batch figure
    (e.g.\ $13639.8\,\text{s} / 48905 = 278.9\,\text{ms}$), so each recorded ``forward call''
    corresponds to one batch of eight circuit executions.
    \item The four-qubit \textbf{Custom~1} configurations achieve a per-sample cost of
    $10.45$\,ms (IndicBERTv2) and $11.07$\,ms (BanglaBERT), which is roughly $3\times$ cheaper
    than the eight-qubit Custom~2 runs ($\approx$35\,ms) and only marginally slower than the
    four-qubit SEL configurations ($8.52$\,ms and $10.46$\,ms respectively). This modest
    overhead relative to SEL is consistent with Custom~1 employing additional single-qubit
    rotation gates per layer while operating on the same $2^4 = 16$-dimensional statevector.
    \item The quantum simulator is the computational bottleneck for the quantum branch, but
    in the full hybrid model it is small relative to the transformer forward pass, which is why
    the total training time is dominated by BERT.
    \item Ranking by per-sample cost at 4 qubits:
    SEL ($\approx$8.5--10.5\,ms) $<$ Custom~1 ($\approx$10.5--11.1\,ms) $\ll$ Custom~2
    at 8 qubits ($\approx$35\,ms), consistent with the exponential $2^{n_q}$ statevector growth
    dominating when the register doubles from 4 to 8 qubits.
\end{itemize}

\section*{3. Single-Label Analysis}

\subsection*{3.1 Dataset Partitioning and Formulation}

In multi-label emotion classification, test samples can be partitioned based on label cardinality into single-label instances ($\sum_{i=1}^L y_i = 1$) representing mutually exclusive emotion states and multi-label co-occurrence instances ($\sum_{i=1}^L y_i > 1$). Filtering the EmoNoBa test set for single-label instances yields $N_{\text{single}} = 940$ test samples across six emotion classes: \emph{Happy} ($N=300, 31.9\%$), \emph{Sad} ($N=200, 21.3\%$), \emph{Angry} ($N=200, 21.3\%$), \emph{Disgust} ($N=100, 10.6\%$), \emph{Surprise} ($N=80, 8.5\%$), and \emph{Fear} ($N=60, 6.4\%$). On this subset, strict Exact Match Ratio (EMR), Micro-F1, Precision, and Recall converge to standard multi-class classification accuracy.

\subsection*{3.2 Single-Label Exact Match Ratio (EMR) Benchmark}

Table~\ref{tab:single_label_emr} summarizes the single-label Exact Match Ratio (EMR) scores across the three investigated VQC notebook implementations: the 4-qubit SEL reference design (\texttt{4qb\_sel.ipynb}), the 4-qubit Custom~1 cross-coupling design (\texttt{Biswanath\_4\_qubit.ipynb}), and the 8-qubit Custom~2 data re-uploading design (\texttt{tathagata\_main\_codebase\_v2 8\_qubit.ipynb}).

\begin{table}[h!]
\centering
\small
\begin{threeparttable}
\begin{tabular}{llcc}
\toprule
\textbf{Encoder} & \textbf{VQC Architecture} & \textbf{Qubits} & \textbf{Exact Match Ratio (EMR)} \\
\midrule
IndicBERTv2 & SEL      & 4 & 0.4713 \\
IndicBERTv2 & Custom 1 & 4 & 0.4521 \\
IndicBERTv2 & Custom 2 & 8 & \textbf{0.5479} \\
\addlinespace
BanglaBERT  & SEL      & 4 & 0.4500 \\
BanglaBERT  & Custom 1 & 4 & 0.3191 \\
BanglaBERT  & Custom 2 & 8 & \textbf{0.5330} \\
\bottomrule
\end{tabular}
\end{threeparttable}
\caption{Single-label Exact Match Ratio (EMR) across evaluated VQC notebook architectures ($N_{\text{single}} = 940$).}
\label{tab:single_label_emr}
\end{table}

\subsection*{3.3 Architectural Investigation and Comparative Analysis}

Analysis of the three VQC notebook implementations reveals how encoding complexity, entanglement topology, and measurement strategies dictate performance on single-label decision boundaries:

\begin{itemize}[leftmargin=*]
    \item \textbf{SEL Architecture (\texttt{4qb\_sel.ipynb}):} Combines a single-pass $R_X(x_i)$ angle embedding with a single layer of PennyLane's \texttt{StronglyEntanglingLayers} (applying general $Rot(\phi,\theta,\omega)$ rotations and default CNOT entanglement) and 4-channel Pauli-$Z$ expectation measurement ($\langle Z_i \rangle$). On single-label instances, SEL achieves moderate EMR ($0.4713$ on IndicBERTv2, $0.4500$ on BanglaBERT) but exhibits extreme majority-class bias, predicting non-zero recall only for high-support categories (\emph{Happy} and \emph{Angry}) while completely failing on minority classes (\emph{Disgust}, \emph{Surprise}, and \emph{Fear} all record $0.0000$ recall).
    
    \item \textbf{Custom 1 Architecture (\texttt{Biswanath\_4\_qubit.ipynb}):} Implements an $R_Y(x_i)$ angle embedding paired with a two-layer ansatz featuring distinct CNOT entangling topologies: ring coupling $(0\to1, 1\to2, 2\to3, 3\to0)$ in Layer 1 and cross-coupling $(0\to2, 1\to3, 2\to1, 3\to0)$ in Layer 2. Despite the deeper CNOT entanglement, the absence of data re-uploading limits feature separability. Under BanglaBERT, Custom~1 experiences complete majority-class collapse, predicting the dominant \emph{Happy} label for all samples (recall $= 1.0000$ on \emph{Happy}, $0.0000$ on all other classes), resulting in an EMR of $0.3191$ (matching the test set proportion $300/940$).
    
    \item \textbf{Custom 2 Architecture (\texttt{tathagata\_main\_codebase\_v2 8\_qubit.ipynb}):} Features a data re-uploading scheme ($U(x) \to V(\theta) \to U(x) \to V(\theta)$) operating on an 8-qubit register. Input features are transformed using learnable per-qubit phase ($\phi_q$) and scale ($\lambda_q$) parameters into paired $R_X$ and $R_Y$ rotations:
    \[
    x_{RX, q} = \lambda_q (\cos\phi_q \cdot x_{2q} - \sin\phi_q \cdot x_{2q+1}), \qquad
    x_{RY, q} = \lambda_q (\sin\phi_q \cdot x_{2q} + \cos\phi_q \cdot x_{2q+1}).
    \]
    The ansatz combines general $Rot$ transformations with ring Controlled-$Z$ ($CZ$) entangling gates ($q \to (q+1)\bmod 8$) and an expanded 16-channel Pauli-$Z$ and Pauli-$X$ readout ($\langle Z_q \rangle, \langle X_q \rangle$). Data re-uploading and phase-scale mixing allow Custom~2 to construct highly non-linear decision boundaries, achieving peak single-label EMR performance ($0.5479$ for IndicBERTv2, $0.5330$ for BanglaBERT) and preventing class-prediction collapse across both language encoders.
\end{itemize}

\end{document}
